\section{核验锚点簇：隆庆、万历：赋役、边市与战争财政（1568--1620）}
\label{register:1568-1620_longqing_wanli}
边市、清丈、漕运和朝鲜战场把资源从不同方向拉向中枢。以下条目不把法定制度等同于地方负担，也不把战报或奏疏中的一时数字推广为全帝国的实况。

以下锚点只承担年序导航；本章前文负责叙事。完整来源、校勘与材料限度留在来源附录和必要的信息框，不把账本字段重复铺成读者段落。

\begin{description}
  \item[\eventanchor{ML-1568-001}{1568（隆庆二年）：俺答边患、东南海防或沿海武装的本年军令与战报构成危机处理线索，战果和参与者须交叉核对。}]
  \item[\eventanchor{ML-1568-002}{1568（隆庆二年）：军饷、盐课、漕运和赋役的本年行政材料显示中晚明财政压力，法定额与实收不可互代。}]
  \item[\eventanchor{ML-1568-003}{1568（隆庆二年）：封贡、互市、月港和港口管理的本年文书记录边海关系重组，官方许可不等于贸易全貌。}]
  \item[\eventanchor{ML-1568-004}{1568（隆庆二年）：海防知识、学校、科举与礼制的本年文本提供制度和知识史线索，方案不等于实际配置。}]
  \item[\eventanchor{MM-1568-QI-JIGUANG-JIZHOU}{1568（隆庆二年）：戚继光调任蓟州总兵，开始整饬蓟镇练兵与防务}]
  \item[\eventanchor{ML-1569-001}{1569（隆庆三年）：军饷、盐课、漕运和赋役的本年行政材料显示中晚明财政压力，法定额与实收不可互代。}]
  \item[\eventanchor{ML-1569-002}{1569（隆庆三年）：封贡、互市、月港和港口管理的本年文书记录边海关系重组，官方许可不等于贸易全貌。}]
  \item[\eventanchor{ML-1569-003}{1569（隆庆三年）：严嵩时期至隆庆初的人事和奏议材料标记中枢权力变化，案狱与党争标签需保留来源立场。}]
  \item[\eventanchor{ML-1569-004}{1569（隆庆三年）：本年灾情、赈恤或河工记录连接中央与地方社会，区域差异需以府县材料追踪。}]
  \item[\eventanchor{ML-1570-001}{1570（隆庆四年）：建州卫王高袭击明朝聚居地}]
  \item[\eventanchor{ML-1570-002}{1570（隆庆四年）：严嵩时期至隆庆初的人事和奏议材料标记中枢权力变化，案狱与党争标签需保留来源立场。}]
  \item[\eventanchor{ML-1570-003}{1570（隆庆四年）：海防知识、学校、科举与礼制的本年文本提供制度和知识史线索，方案不等于实际配置。}]
  \item[\eventanchor{ML-1570-004}{1570（隆庆四年）：俺答边患、东南海防或沿海武装的本年军令与战报构成危机处理线索，战果和参与者须交叉核对。}]
  \item[\eventanchor{MM-1570-ALTAN-NEGOTIATIONS}{1570（隆庆四年）：明廷与俺答部围绕封贡互市及赵全等人问题展开谈判}]
  \item[\eventanchor{ML-1571-001}{1571（隆庆五年）：封贡、互市、月港和港口管理的本年文书记录边海关系重组，官方许可不等于贸易全貌。}]
  \item[\eventanchor{ML-1571-002}{1571（隆庆五年）：海防知识、学校、科举与礼制的本年文本提供制度和知识史线索，方案不等于实际配置。}]
  \item[\eventanchor{ML-1571-003}{1571（隆庆五年）：本年灾情、赈恤或河工记录连接中央与地方社会，区域差异需以府县材料追踪。}]
  \item[\eventanchor{ML-1571-004}{1571（隆庆五年）：军饷、盐课、漕运和赋役的本年行政材料显示中晚明财政压力，法定额与实收不可互代。}]
  \item[\eventanchor{ML-1572-001}{1572（隆庆六年）七月5日：隆庆皇帝驾崩}]
  \item[\eventanchor{ML-1572-002}{1572（隆庆六年）七月19日：朱翊钧成为万历皇帝}]
  \item[\eventanchor{ML-1572-003}{1572（隆庆六年）：封贡、互市、月港和港口管理的本年文书记录边海关系重组，官方许可不等于贸易全貌。}]
  \item[\eventanchor{MM-1572-GAO-GONG-DISMISSAL}{1572（隆庆六年）六月：高拱被罢，张居正成为万历初年内阁的核心人物}]
  \item[\eventanchor{MM-1572-WANLI-ACCESSION}{1572（隆庆六年）五月：隆庆帝去世，十岁的朱翊钧继位，是为万历帝}]
  \item[\eventanchor{ML-1573-001}{1573（万历元年）：西班牙与中国的贸易从月港开始}]
  \item[\eventanchor{ML-1573-002}{1573（万历元年）：俺答边患、东南海防或沿海武装的本年军令与战报构成危机处理线索，战果和参与者须交叉核对。}]
  \item[\eventanchor{ML-1573-003}{1573（万历元年）：严嵩时期至隆庆初的人事和奏议材料标记中枢权力变化，案狱与党争标签需保留来源立场。}]
  \item[\eventanchor{ML-1574-001}{1574（万历二年）：李成梁在乔仓嘎和塔克西的帮助下杀死了王高}]
  \item[\eventanchor{ML-1574-002}{1574（万历二年）：澳门周围竖起了围墙}]
  \item[\eventanchor{ML-1574-003}{1574（万历二年）：本年灾情、赈恤或河工记录连接中央与地方社会，区域差异需以府县材料追踪。}]
  \item[\eventanchor{ML-1574-004}{1574（万历二年）：军饷、盐课、漕运和赋役的本年行政材料显示中晚明财政压力，法定额与实收不可互代。}]
  \item[\eventanchor{ML-1574-005}{1574（万历二年）：海防知识、学校、科举与礼制的本年文本提供制度和知识史线索，方案不等于实际配置。}]
  \item[\eventanchor{ML-1575-001}{1575（万历三年）：明朝海军军官王旺高抵达吕宋岛，并随马丁·德·拉达率领的西班牙大使馆返回；由于西班牙无法抓获中国海盗林峰，使馆失败}]
  \item[\eventanchor{ML-1575-002}{1575（万历三年）：考成、人事、立储和留中等本年材料标记皇帝、内阁与言官的协调关系，不可只用党争标签解释。（材料核查重点：诏令或奏议的形成与发受链）}]
  \item[\eventanchor{ML-1575-003}{1575（万历三年）：清丈、一条鞭、漕运或军饷的本年文书构成万历财政的可核查节点，定额、报数和实收必须分列。}]
  \item[\eventanchor{ML-1575-004}{1575（万历三年）：考成、人事、立储和留中等本年材料标记皇帝、内阁与言官的协调关系，不可只用党争标签解释。（材料核查重点：报称信息的来源、抄传与行政分类）}]
  \item[\eventanchor{ML-1575-005}{1575（万历三年）：书院、科举、出版与宗教活动的本年线索可连接知识网络，存世文本有阶层与地域偏差。}]
  \item[\eventanchor{ML-1575-006}{1575（万历三年）：辽东、西北、朝鲜或西南军务的本年调兵与战报记录跨区域动员，补给与兵额需要独立核验。}]
  \item[\eventanchor{ML-1576-001}{1576（万历四年）：中美贸易建立}]
  \item[\eventanchor{ML-1576-002}{1576（万历四年）：辽东、西北、朝鲜或西南军务的本年调兵与战报记录跨区域动员，补给与兵额需要独立核验。（材料核查重点：诏令或奏议的形成与发受链）}]
  \item[\eventanchor{ML-1576-003}{1576（万历四年）：考成、人事、立储和留中等本年材料标记皇帝、内阁与言官的协调关系，不可只用党争标签解释。}]
  \item[\eventanchor{ML-1576-004}{1576（万历四年）：辽东、西北、朝鲜或西南军务的本年调兵与战报记录跨区域动员，补给与兵额需要独立核验。（材料核查重点：报称信息的来源、抄传与行政分类）}]
  \item[\eventanchor{ML-1576-005}{1576（万历四年）：本年灾荒、河工和地方赈济材料可用于研究风险分配，灾害叙事和统计数字应区分。}]
  \item[\eventanchor{ML-1576-006}{1576（万历四年）：朝鲜战争、海上贸易和朝贡使节的本年多方记录提供跨政权比较，外交语言不能替代地方经验。}]
  \item[\eventanchor{ML-1577-001}{1577（万历五年）：辽东、西北、朝鲜或西南军务的本年调兵与战报记录跨区域动员，补给与兵额需要独立核验。}]
  \item[\eventanchor{ML-1577-003}{1577（万历五年）：清丈、一条鞭、漕运或军饷的本年文书构成万历财政的可核查节点，定额、报数和实收必须分列。}]
  \item[\eventanchor{ML-1577-004}{1577（万历五年）：书院、科举、出版与宗教活动的本年线索可连接知识网络，存世文本有阶层与地域偏差。}]
  \item[\eventanchor{MM-1577-ZHANG-JUZHENG-DUOQING}{1577（万历五年）：张居正父丧期间奉旨夺情留任，引发朝臣围绕礼制与政务的争论}]
  \item[\eventanchor{ML-1578-001}{1578（万历六年）：葡萄牙人获准前往广州}]
  \item[\eventanchor{ML-1578-002}{1578（万历六年）：书院、科举、出版与宗教活动的本年线索可连接知识网络，存世文本有阶层与地域偏差。}]
  \item[\eventanchor{ML-1578-003}{1578（万历六年）：考成、人事、立储和留中等本年材料标记皇帝、内阁与言官的协调关系，不可只用党争标签解释。}]
  \item[\eventanchor{ML-1578-004}{1578（万历六年）：本年灾荒、河工和地方赈济材料可用于研究风险分配，灾害叙事和统计数字应区分。}]
  \item[\eventanchor{MM-1578-EMPRESS-WANG-INSTALLATION}{1578（万历六年）六月：万历帝册立王氏为皇后，完成幼帝成年后的皇后册立仪式}]
  \item[\eventanchor{ML-1579-001}{1579（万历七年）：东林运动：私立东林书院全部关闭}]
  \item[\eventanchor{ML-1579-002}{1579（万历七年）：朝鲜战争、海上贸易和朝贡使节的本年多方记录提供跨政权比较，外交语言不能替代地方经验。}]
  \item[\eventanchor{ML-1579-003}{1579（万历七年）：本年灾荒、河工和地方赈济材料可用于研究风险分配，灾害叙事和统计数字应区分。}]
  \item[\eventanchor{ML-1579-004}{1579（万历七年）：辽东、西北、朝鲜或西南军务的本年调兵与战报记录跨区域动员，补给与兵额需要独立核验。}]
  \item[\eventanchor{ML-1579-005}{1579（万历七年）：清丈、一条鞭、漕运或军饷的本年文书构成万历财政的可核查节点，定额、报数和实收必须分列。}]
  \item[\eventanchor{ML-1580-001}{1580（万历八年）：单鞭法：简化税法}]
  \item[\eventanchor{ML-1580-002}{1580（万历八年）：官员们批评万历皇帝的疏忽和个人生活的不当行为}]
  \item[\eventanchor{ML-1580-003}{1580（万历八年）：清丈、一条鞭、漕运或军饷的本年文书构成万历财政的可核查节点，定额、报数和实收必须分列。}]
  \item[\eventanchor{ML-1580-004}{1580（万历八年）：朝鲜战争、海上贸易和朝贡使节的本年多方记录提供跨政权比较，外交语言不能替代地方经验。}]
  \item[\eventanchor{ML-1581-001}{1581（万历九年）：朝鲜战争、海上贸易和朝贡使节的本年多方记录提供跨政权比较，外交语言不能替代地方经验。}]
  \item[\eventanchor{ML-1581-002}{1581（万历九年）：书院、科举、出版与宗教活动的本年线索可连接知识网络，存世文本有阶层与地域偏差。（材料核查重点：诏令或奏议的形成与发受链）}]
  \item[\eventanchor{ML-1581-003}{1581（万历九年）：考成、人事、立储和留中等本年材料标记皇帝、内阁与言官的协调关系，不可只用党争标签解释。}]
  \item[\eventanchor{ML-1581-004}{1581（万历九年）：书院、科举、出版与宗教活动的本年线索可连接知识网络，存世文本有阶层与地域偏差。（材料核查重点：报称信息的来源、抄传与行政分类）}]
  \item[\eventanchor{ML-1581-005}{1581（万历九年）：辽东、西北、朝鲜或西南军务的本年调兵与战报记录跨区域动员，补给与兵额需要独立核验。}]
  \item[\eventanchor{ML-1582-001}{1582（万历十年）：明军击败建州女真阿泰，误杀努尔哈赤的祖父和父亲乔仓嘎和塔克西}]
  \item[\eventanchor{ML-1582-002}{1582（万历十年）：太仓国库积存白银超600万两}]
  \item[\eventanchor{ML-1582-003}{1582（万历十年）：欧洲生产比中国历更准确的公历}]
  \item[\eventanchor{ML-1582-004}{1582（万历十年）：本年灾荒、河工和地方赈济材料可用于研究风险分配，灾害叙事和统计数字应区分。}]
  \item[\eventanchor{ML-1582-005}{1582（万历十年）：朝鲜战争、海上贸易和朝贡使节的本年多方记录提供跨政权比较，外交语言不能替代地方经验。}]
  \item[\eventanchor{ML-1583-001}{1583（万历十一年）：利玛窦在肇庆设立教堂}]
  \item[\eventanchor{ML-1583-002}{1583（万历十一年）：书院、科举、出版与宗教活动的本年线索可连接知识网络，存世文本有阶层与地域偏差。（材料核查重点：诏令或奏议的形成与发受链）}]
  \item[\eventanchor{ML-1583-003}{1583（万历十一年）：本年灾荒、河工和地方赈济材料可用于研究风险分配，灾害叙事和统计数字应区分。}]
  \item[\eventanchor{ML-1583-004}{1583（万历十一年）：辽东、西北、朝鲜或西南军务的本年调兵与战报记录跨区域动员，补给与兵额需要独立核验。}]
  \item[\eventanchor{ML-1583-005}{1583（万历十一年）：清丈、一条鞭、漕运或军饷的本年文书构成万历财政的可核查节点，定额、报数和实收必须分列。}]
  \item[\eventanchor{ML-1583-006}{1583（万历十一年）：书院、科举、出版与宗教活动的本年线索可连接知识网络，存世文本有阶层与地域偏差。（材料核查重点：报称信息的来源、抄传与行政分类）}]
  \item[\eventanchor{ML-1584-001}{1584（万历十二年）：清丈、一条鞭、漕运或军饷的本年文书构成万历财政的可核查节点，定额、报数和实收必须分列。（材料核查重点：诏令或奏议的形成与发受链）}]
  \item[\eventanchor{ML-1584-002}{1584（万历十二年）：本年灾荒、河工和地方赈济材料可用于研究风险分配，灾害叙事和统计数字应区分。（材料核查重点：诏令或奏议的形成与发受链）}]
  \item[\eventanchor{ML-1584-004}{1584（万历十二年）：朝鲜战争、海上贸易和朝贡使节的本年多方记录提供跨政权比较，外交语言不能替代地方经验。}]
  \item[\eventanchor{ML-1584-005}{1584（万历十二年）：考成、人事、立储和留中等本年材料标记皇帝、内阁与言官的协调关系，不可只用党争标签解释。}]
  \item[\eventanchor{ML-1584-006}{1584（万历十二年）：本年灾荒、河工和地方赈济材料可用于研究风险分配，灾害叙事和统计数字应区分。（材料核查重点：报称信息的来源、抄传与行政分类）}]
  \item[\eventanchor{ML-1585-001}{1585（万历十三年）：考成、人事、立储和留中等本年材料标记皇帝、内阁与言官的协调关系，不可只用党争标签解释。（材料核查重点：诏令或奏议的形成与发受链）}]
  \item[\eventanchor{ML-1585-002}{1585（万历十三年）：清丈、一条鞭、漕运或军饷的本年文书构成万历财政的可核查节点，定额、报数和实收必须分列。（材料核查重点：诏令或奏议的形成与发受链）}]
  \item[\eventanchor{ML-1585-003}{1585（万历十三年）：考成、人事、立储和留中等本年材料标记皇帝、内阁与言官的协调关系，不可只用党争标签解释。（材料核查重点：报称信息的来源、抄传与行政分类）}]
  \item[\eventanchor{ML-1585-004}{1585（万历十三年）：书院、科举、出版与宗教活动的本年线索可连接知识网络，存世文本有阶层与地域偏差。}]
  \item[\eventanchor{ML-1585-005}{1585（万历十三年）：辽东、西北、朝鲜或西南军务的本年调兵与战报记录跨区域动员，补给与兵额需要独立核验。}]
  \item[\eventanchor{ML-1585-006}{1585（万历十三年）：清丈、一条鞭、漕运或军饷的本年文书构成万历财政的可核查节点，定额、报数和实收必须分列。（材料核查重点：报称信息的来源、抄传与行政分类）}]
  \item[\eventanchor{ML-1586-001}{1586（万历十四年）：辽东、西北、朝鲜或西南军务的本年调兵与战报记录跨区域动员，补给与兵额需要独立核验。（材料核查重点：诏令或奏议的形成与发受链）}]
  \item[\eventanchor{ML-1586-002}{1586（万历十四年）：考成、人事、立储和留中等本年材料标记皇帝、内阁与言官的协调关系，不可只用党争标签解释。（材料核查重点：诏令或奏议的形成与发受链）}]
  \item[\eventanchor{ML-1586-003}{1586（万历十四年）：辽东、西北、朝鲜或西南军务的本年调兵与战报记录跨区域动员，补给与兵额需要独立核验。（材料核查重点：报称信息的来源、抄传与行政分类）}]
  \item[\eventanchor{ML-1586-004}{1586（万历十四年）：本年灾荒、河工和地方赈济材料可用于研究风险分配，灾害叙事和统计数字应区分。}]
  \item[\eventanchor{ML-1586-005}{1586（万历十四年）：朝鲜战争、海上贸易和朝贡使节的本年多方记录提供跨政权比较，外交语言不能替代地方经验。}]
  \item[\eventanchor{ML-1586-006}{1586（万历十四年）：考成、人事、立储和留中等本年材料标记皇帝、内阁与言官的协调关系，不可只用党争标签解释。（材料核查重点：报称信息的来源、抄传与行政分类）}]
  \item[\eventanchor{ML-1587-001}{1587（万历十五年）：朝鲜战争、海上贸易和朝贡使节的本年多方记录提供跨政权比较，外交语言不能替代地方经验。（材料核查重点：诏令或奏议的形成与发受链）}]
  \item[\eventanchor{ML-1587-002}{1587（万历十五年）：辽东、西北、朝鲜或西南军务的本年调兵与战报记录跨区域动员，补给与兵额需要独立核验。（材料核查重点：诏令或奏议的形成与发受链）}]
  \item[\eventanchor{ML-1587-004}{1587（万历十五年）：清丈、一条鞭、漕运或军饷的本年文书构成万历财政的可核查节点，定额、报数和实收必须分列。}]
  \item[\eventanchor{ML-1587-005}{1587（万历十五年）：书院、科举、出版与宗教活动的本年线索可连接知识网络，存世文本有阶层与地域偏差。}]
  \item[\eventanchor{ML-1587-006}{1587（万历十五年）：辽东、西北、朝鲜或西南军务的本年调兵与战报记录跨区域动员，补给与兵额需要独立核验。（材料核查重点：报称信息的来源、抄传与行政分类）}]
  \item[\eventanchor{ML-1588-001}{1588（万历十六年）：书院、科举、出版与宗教活动的本年线索可连接知识网络，存世文本有阶层与地域偏差。（材料核查重点：诏令或奏议的形成与发受链）}]
  \item[\eventanchor{ML-1588-002}{1588（万历十六年）：朝鲜战争、海上贸易和朝贡使节的本年多方记录提供跨政权比较，外交语言不能替代地方经验。}]
  \item[\eventanchor{ML-1588-003}{1588（万历十六年）：书院、科举、出版与宗教活动的本年线索可连接知识网络，存世文本有阶层与地域偏差。（材料核查重点：报称信息的来源、抄传与行政分类）}]
  \item[\eventanchor{ML-1588-004}{1588（万历十六年）：考成、人事、立储和留中等本年材料标记皇帝、内阁与言官的协调关系，不可只用党争标签解释。}]
  \item[\eventanchor{ML-1588-005}{1588（万历十六年）：本年灾荒、河工和地方赈济材料可用于研究风险分配，灾害叙事和统计数字应区分。}]
  \item[\eventanchor{ML-1589-001}{1589（万历十七年）：亳州叛乱：亳州苗族叛乱}]
  \item[\eventanchor{ML-1589-002}{1589（万历十七年）：书院、科举、出版与宗教活动的本年线索可连接知识网络，存世文本有阶层与地域偏差。}]
  \item[\eventanchor{ML-1589-003}{1589（万历十七年）：本年灾荒、河工和地方赈济材料可用于研究风险分配，灾害叙事和统计数字应区分。}]
  \item[\eventanchor{ML-1589-004}{1589（万历十七年）：辽东、西北、朝鲜或西南军务的本年调兵与战报记录跨区域动员，补给与兵额需要独立核验。}]
  \item[\eventanchor{ML-1589-005}{1589（万历十七年）：清丈、一条鞭、漕运或军饷的本年文书构成万历财政的可核查节点，定额、报数和实收必须分列。}]
  \item[\eventanchor{ML-1590-001}{1590（万历十八年）：福建华人开始在台湾西南部定居}]
  \item[\eventanchor{ML-1590-002}{1590（万历十八年）：本年灾荒、河工和地方赈济材料可用于研究风险分配，灾害叙事和统计数字应区分。}]
  \item[\eventanchor{ML-1590-003}{1590（万历十八年）：清丈、一条鞭、漕运或军饷的本年文书构成万历财政的可核查节点，定额、报数和实收必须分列。}]
  \item[\eventanchor{ML-1590-004}{1590（万历十八年）：朝鲜战争、海上贸易和朝贡使节的本年多方记录提供跨政权比较，外交语言不能替代地方经验。}]
  \item[\eventanchor{ML-1590-005}{1590（万历十八年）：考成、人事、立储和留中等本年材料标记皇帝、内阁与言官的协调关系，不可只用党争标签解释。}]
  \item[\eventanchor{ML-1591-002}{1591（万历十九年）：清丈、一条鞭、漕运或军饷的本年文书构成万历财政的可核查节点，定额、报数和实收必须分列。}]
  \item[\eventanchor{ML-1591-003}{1591（万历十九年）：考成、人事、立储和留中等本年材料标记皇帝、内阁与言官的协调关系，不可只用党争标签解释。}]
  \item[\eventanchor{ML-1591-004}{1591（万历十九年）：书院、科举、出版与宗教活动的本年线索可连接知识网络，存世文本有阶层与地域偏差。}]
  \item[\eventanchor{ML-1591-005}{1591（万历十九年）：辽东、西北、朝鲜或西南军务的本年调兵与战报记录跨区域动员，补给与兵额需要独立核验。}]
  \item[\eventanchor{ML-1592-001}{1592（万历二十年）三月：鄂尔多斯战役：宁夏刘东阳和浦北叛军}]
  \item[\eventanchor{ML-1592-002}{1592（万历二十年）七月14日：鄂尔多斯战役：叶梦雄携大炮增援苗军围攻宁夏}]
  \item[\eventanchor{ML-1592-003}{1592（万历二十年）八月23日：鄂尔多斯战役：宁夏周边堤坝竣工}]
  \item[\eventanchor{ML-1592-004}{1592（万历二十年）九月6日：鄂尔多斯战役：宁夏被淹}]
  \item[\eventanchor{ML-1592-005}{1592（万历二十年）九月25日：鄂尔多斯战役：叛军最后一次尝试突围宁夏}]
  \item[\eventanchor{ML-1592-006}{1592（万历二十年）十月12日：鄂尔多斯战役：北墙倒塌，叛乱失败}]
  \item[\eventanchor{ML-1592-007}{1592（万历二十年）：日本入侵朝鲜（1592-98）：明朝军队动员干预日本入侵朝鲜}]
  \item[\eventanchor{ML-1592-008}{1592（万历二十年）：日本入侵朝鲜（1592-98）：努尔哈赤提出与日本作战，但遭到拒绝；明朝对努尔哈赤军队的规模和素质感到震惊}]
  \item[\eventanchor{ML-1593-001}{1593（万历二十一年）一月8日：平壤之围（1593 年）：明朝军队将日本军队逐出平壤}]
  \item[\eventanchor{ML-1593-002}{1593（万历二十一年）一月27日：碧蹄馆之战：明朝先头侦察队被日军击败}]
  \item[\eventanchor{ML-1593-003}{1593（万历二十一年）：嘉运河扩建开始}]
  \item[\eventanchor{ML-1593-004}{1593（万历二十一年）：中、初年级作业采用抽签方式分配}]
  \item[\eventanchor{ML-1593-005}{1593（万历二十一年）：明朝官员每年向中国帆船颁发十张许可证，允许其在台湾北部进行贸易}]
  \item[\eventanchor{ML-1594-002}{1594（万历二十二年）：考成、人事、立储和留中等本年材料标记皇帝、内阁与言官的协调关系，不可只用党争标签解释。}]
  \item[\eventanchor{ML-1594-003}{1594（万历二十二年）：辽东、西北、朝鲜或西南军务的本年调兵与战报记录跨区域动员，补给与兵额需要独立核验。}]
  \item[\eventanchor{ML-1594-004}{1594（万历二十二年）：本年灾荒、河工和地方赈济材料可用于研究风险分配，灾害叙事和统计数字应区分。}]
  \item[\eventanchor{ML-1594-005}{1594（万历二十二年）：朝鲜战争、海上贸易和朝贡使节的本年多方记录提供跨政权比较，外交语言不能替代地方经验。}]
  \item[\eventanchor{ML-1595-001}{1595（万历二十三年）：清丈、一条鞭、漕运或军饷的本年文书构成万历财政的可核查节点，定额、报数和实收必须分列。（材料核查重点：诏令或奏议的形成与发受链）}]
  \item[\eventanchor{ML-1595-002}{1595（万历二十三年）：辽东、西北、朝鲜或西南军务的本年调兵与战报记录跨区域动员，补给与兵额需要独立核验。}]
  \item[\eventanchor{ML-1595-003}{1595（万历二十三年）：朝鲜战争、海上贸易和朝贡使节的本年多方记录提供跨政权比较，外交语言不能替代地方经验。}]
  \item[\eventanchor{ML-1595-004}{1595（万历二十三年）：清丈、一条鞭、漕运或军饷的本年文书构成万历财政的可核查节点，定额、报数和实收必须分列。（材料核查重点：报称信息的来源、抄传与行政分类）}]
  \item[\eventanchor{ML-1595-005}{1595（万历二十三年）：书院、科举、出版与宗教活动的本年线索可连接知识网络，存世文本有阶层与地域偏差。}]
  \item[\eventanchor{ML-1596-001}{1596（万历二十四年）：万历皇帝派遣宦官担任税吏和矿监}]
  \item[\eventanchor{ML-1596-002}{1596（万历二十四年）：朝鲜战争、海上贸易和朝贡使节的本年多方记录提供跨政权比较，外交语言不能替代地方经验。}]
  \item[\eventanchor{ML-1596-003}{1596（万历二十四年）：书院、科举、出版与宗教活动的本年线索可连接知识网络，存世文本有阶层与地域偏差。}]
  \item[\eventanchor{ML-1596-004}{1596（万历二十四年）：考成、人事、立储和留中等本年材料标记皇帝、内阁与言官的协调关系，不可只用党争标签解释。}]
  \item[\eventanchor{ML-1596-005}{1596（万历二十四年）：本年灾荒、河工和地方赈济材料可用于研究风险分配，灾害叙事和统计数字应区分。}]
  \item[\eventanchor{ML-1597-001}{1597（万历二十五年）十月17日：稷山之战：日军向汉城进军被明军阻止}]
  \item[\eventanchor{ML-1597-002}{1597（万历二十五年）：书院、科举、出版与宗教活动的本年线索可连接知识网络，存世文本有阶层与地域偏差。}]
  \item[\eventanchor{ML-1597-003}{1597（万历二十五年）：本年灾荒、河工和地方赈济材料可用于研究风险分配，灾害叙事和统计数字应区分。}]
  \item[\eventanchor{ML-1597-004}{1597（万历二十五年）：辽东、西北、朝鲜或西南军务的本年调兵与战报记录跨区域动员，补给与兵额需要独立核验。}]
  \item[\eventanchor{ML-1598-001}{1598（万历二十六年）一月4日：蔚山之围：明朝和朝鲜军队未能将日本人赶出蔚山城堡}]
  \item[\eventanchor{ML-1598-002}{1598（万历二十六年）十月：泗川之战（1598）：明朝和朝鲜军队未能将日本人赶出泗川}]
  \item[\eventanchor{ML-1598-003}{1598（万历二十六年）十月7日：围攻顺天：明朝和朝鲜军队未能将日本人逐出顺天城}]
  \item[\eventanchor{ML-1598-004}{1598（万历二十六年）十二月16日：鹭梁海战：明朝鲜海军击败日本舰队}]
  \item[\eventanchor{ML-1598-005}{1598（万历二十六年）十二月24日：日本入侵朝鲜（1592-98）：日本军队从朝鲜撤军}]
  \item[\eventanchor{ML-1598-006}{1598（万历二十六年）：亳州叛乱：苗族叛乱被镇压}]
  \item[\eventanchor{ML-1598-007}{1598（万历二十六年）：蒙古人杀明主帅李如松}]
  \item[\eventanchor{ML-1598-008}{1598（万历二十六年）：明朝骑兵尝试发射三管火绳枪，然后用它作为盾牌，同时用另一只手用马刀攻击。}]
  \item[\eventanchor{ML-1598-009}{1598（万历二十六年）：广东官员允许西班牙人在皮纳尔进行贸易}]
  \item[\eventanchor{ML-1599-001}{1599（万历二十七年）：各大港口均有高级太监驻守}]
  \item[\eventanchor{ML-1599-002}{1599（万历二十七年）：本年灾荒、河工和地方赈济材料可用于研究风险分配，灾害叙事和统计数字应区分。}]
  \item[\eventanchor{ML-1599-003}{1599（万历二十七年）：清丈、一条鞭、漕运或军饷的本年文书构成万历财政的可核查节点，定额、报数和实收必须分列。}]
  \item[\eventanchor{ML-1599-004}{1599（万历二十七年）：朝鲜战争、海上贸易和朝贡使节的本年多方记录提供跨政权比较，外交语言不能替代地方经验。}]
  \item[\eventanchor{ML-1599-005}{1599（万历二十七年）：考成、人事、立储和留中等本年材料标记皇帝、内阁与言官的协调关系，不可只用党争标签解释。}]
  \item[\eventanchor{ML-1600-001}{1600（万历二十八年）一月17：澳门的葡萄牙人攻击兰帕考的西班牙人。西班牙人放弃埃尔皮纳尔。}]
  \item[\eventanchor{ML-1600-002}{1600（万历二十八年）：欧洲藏书规模超过中国}]
  \item[\eventanchor{ML-1600-003}{1600（万历二十八年）：东林、书院、刻书和宗教活动的本年材料记录精英网络，社团存在不等于一致政治立场。}]
  \item[\eventanchor{ML-1600-005}{1600（万历二十八年）：后金兴起、朝鲜交涉和海上网络的本年多方材料揭示区域关系变化，政权名称与译名须按原文说明。}]
  \item[\eventanchor{ML-1600-006}{1600（万历二十八年）：辽东战事、边镇防务和西南兵事的本年军令记录资源调动，战报数字应与朝鲜、满文或地方材料比照。}]
  \item[\eventanchor{ML-1601-001}{1601（万历二十九年）：万历末至天启初的立储、内阁和言路材料标记中枢运作的堵塞与调整，派系归类需回到具体文书。（材料核查重点：诏令或奏议的形成与发受链）}]
  \item[\eventanchor{ML-1601-002}{1601（万历二十九年）：矿税、加派、漕运和辽饷的本年奏疏与账目提供财政压力线索，预算、欠额和实解不得混同。}]
  \item[\eventanchor{ML-1601-003}{1601（万历二十九年）：本年灾害、疫病或地方赈济线索应以受灾地区文书和地方志复核，中央奏报不能覆盖全部社会。}]
  \item[\eventanchor{ML-1601-004}{1601（万历二十九年）：万历末至天启初的立储、内阁和言路材料标记中枢运作的堵塞与调整，派系归类需回到具体文书。（材料核查重点：报称信息的来源、抄传与行政分类）}]
  \item[\eventanchor{ML-1601-005}{1601（万历二十九年）：东林、书院、刻书和宗教活动的本年材料记录精英网络，社团存在不等于一致政治立场。}]
  \item[\eventanchor{ML-1601-006}{1601（万历二十九年）：后金兴起、朝鲜交涉和海上网络的本年多方材料揭示区域关系变化，政权名称与译名须按原文说明。}]
  \item[\eventanchor{ML-1602-001}{1602（万历三十年）：利玛窦定居北京传播基督教}]
  \item[\eventanchor{ML-1602-002}{1602（万历三十年）：辽东战事、边镇防务和西南兵事的本年军令记录资源调动，战报数字应与朝鲜、满文或地方材料比照。}]
  \item[\eventanchor{ML-1602-003}{1602（万历三十年）：万历末至天启初的立储、内阁和言路材料标记中枢运作的堵塞与调整，派系归类需回到具体文书。}]
  \item[\eventanchor{ML-1602-004}{1602（万历三十年）：矿税、加派、漕运和辽饷的本年奏疏与账目提供财政压力线索，预算、欠额和实解不得混同。}]
  \item[\eventanchor{ML-1602-005}{1602（万历三十年）：本年灾害、疫病或地方赈济线索应以受灾地区文书和地方志复核，中央奏报不能覆盖全部社会。}]
  \item[\eventanchor{ML-1602-006}{1602（万历三十年）：东林、书院、刻书和宗教活动的本年材料记录精英网络，社团存在不等于一致政治立场。}]
  \item[\eventanchor{ML-1603-001}{1603（万历三十一年）十月：桑利叛乱：西班牙人、日本人和菲律宾人屠杀马尼拉的华人；万历皇帝指责一名太监询问西班牙人是否可以在甲美地采矿，激怒了西班牙人}]
  \item[\eventanchor{ML-1603-002}{1603（万历三十一年）：努尔哈赤与明朝将军同意划定领土边界}]
  \item[\eventanchor{ML-1603-003}{1603（万历三十一年）：中国学者陈迪在明代反海盗任务期间在大佑安湾（台湾由此得名）呆了一段时间，首次对台湾原住民进行了重要描述}]
  \item[\eventanchor{ML-1603-004}{1603（万历三十一年）：辽东战事、边镇防务和西南兵事的本年军令记录资源调动，战报数字应与朝鲜、满文或地方材料比照。}]
  \item[\eventanchor{ML-1603-006}{1603（万历三十一年）：本年灾害、疫病或地方赈济线索应以受灾地区文书和地方志复核，中央奏报不能覆盖全部社会。}]
  \item[\eventanchor{ML-1604-001}{1604（万历三十二年）：东林运动：东林书院成立}]
  \item[\eventanchor{ML-1604-002}{1604（万历三十二年）：后金兴起、朝鲜交涉和海上网络的本年多方材料揭示区域关系变化，政权名称与译名须按原文说明。（材料核查重点：诏令或奏议的形成与发受链）}]
  \item[\eventanchor{ML-1604-003}{1604（万历三十二年）：矿税、加派、漕运和辽饷的本年奏疏与账目提供财政压力线索，预算、欠额和实解不得混同。（材料核查重点：诏令或奏议的形成与发受链）}]
  \item[\eventanchor{ML-1604-004}{1604（万历三十二年）：后金兴起、朝鲜交涉和海上网络的本年多方材料揭示区域关系变化，政权名称与译名须按原文说明。（材料核查重点：报称信息的来源、抄传与行政分类）}]
  \item[\eventanchor{ML-1604-005}{1604（万历三十二年）：矿税、加派、漕运和辽饷的本年奏疏与账目提供财政压力线索，预算、欠额和实解不得混同。（材料核查重点：报称信息的来源、抄传与行政分类）}]
  \item[\eventanchor{ML-1604-006}{1604（万历三十二年）：万历末至天启初的立储、内阁和言路材料标记中枢运作的堵塞与调整，派系归类需回到具体文书。}]
  \item[\eventanchor{ML-1605-001}{1605（万历三十三年）六月：一道雷电击倒了天坛的旗杆，非常不吉利，导致一些官员辞职}]
  \item[\eventanchor{ML-1605-002}{1605（万历三十三年）：东林、书院、刻书和宗教活动的本年材料记录精英网络，社团存在不等于一致政治立场。（材料核查重点：诏令或奏议的形成与发受链）}]
  \item[\eventanchor{ML-1605-003}{1605（万历三十三年）：辽东战事、边镇防务和西南兵事的本年军令记录资源调动，战报数字应与朝鲜、满文或地方材料比照。（材料核查重点：诏令或奏议的形成与发受链）}]
  \item[\eventanchor{ML-1605-004}{1605（万历三十三年）：东林、书院、刻书和宗教活动的本年材料记录精英网络，社团存在不等于一致政治立场。（材料核查重点：报称信息的来源、抄传与行政分类）}]
  \item[\eventanchor{ML-1605-005}{1605（万历三十三年）：辽东战事、边镇防务和西南兵事的本年军令记录资源调动，战报数字应与朝鲜、满文或地方材料比照。（材料核查重点：报称信息的来源、抄传与行政分类）}]
  \item[\eventanchor{ML-1605-006}{1605（万历三十三年）：矿税、加派、漕运和辽饷的本年奏疏与账目提供财政压力线索，预算、欠额和实解不得混同。}]
  \item[\eventanchor{ML-1606-001}{1606（万历三十四年）：云南军官骚乱并杀害采矿太监杨荣}]
  \item[\eventanchor{ML-1606-002}{1606（万历三十四年）：明式步枪均附有插入式刺刀。}]
  \item[\eventanchor{ML-1606-003}{1606（万历三十四年）：后金兴起、朝鲜交涉和海上网络的本年多方材料揭示区域关系变化，政权名称与译名须按原文说明。（材料核查重点：诏令或奏议的形成与发受链）}]
  \item[\eventanchor{ML-1606-004}{1606（万历三十四年）：本年灾害、疫病或地方赈济线索应以受灾地区文书和地方志复核，中央奏报不能覆盖全部社会。}]
  \item[\eventanchor{ML-1606-006}{1606（万历三十四年）：辽东战事、边镇防务和西南兵事的本年军令记录资源调动，战报数字应与朝鲜、满文或地方材料比照。}]
  \item[\eventanchor{ML-1607-001}{1607（万历三十五年）：欧几里得几何原理前六册被翻译成中文}]
  \item[\eventanchor{ML-1607-002}{1607（万历三十五年）：本年灾害、疫病或地方赈济线索应以受灾地区文书和地方志复核，中央奏报不能覆盖全部社会。}]
  \item[\eventanchor{ML-1607-003}{1607（万历三十五年）：东林、书院、刻书和宗教活动的本年材料记录精英网络，社团存在不等于一致政治立场。（材料核查重点：诏令或奏议的形成与发受链）}]
  \item[\eventanchor{ML-1607-004}{1607（万历三十五年）：万历末至天启初的立储、内阁和言路材料标记中枢运作的堵塞与调整，派系归类需回到具体文书。}]
  \item[\eventanchor{ML-1607-005}{1607（万历三十五年）：东林、书院、刻书和宗教活动的本年材料记录精英网络，社团存在不等于一致政治立场。（材料核查重点：报称信息的来源、抄传与行政分类）}]
  \item[\eventanchor{ML-1607-006}{1607（万历三十五年）：后金兴起、朝鲜交涉和海上网络的本年多方材料揭示区域关系变化，政权名称与译名须按原文说明。}]
  \item[\eventanchor{ML-1608-001}{1608（万历三十六年）：矿税、加派、漕运和辽饷的本年奏疏与账目提供财政压力线索，预算、欠额和实解不得混同。（材料核查重点：诏令或奏议的形成与发受链）}]
  \item[\eventanchor{ML-1608-002}{1608（万历三十六年）：万历末至天启初的立储、内阁和言路材料标记中枢运作的堵塞与调整，派系归类需回到具体文书。}]
  \item[\eventanchor{ML-1608-003}{1608（万历三十六年）：本年灾害、疫病或地方赈济线索应以受灾地区文书和地方志复核，中央奏报不能覆盖全部社会。（材料核查重点：诏令或奏议的形成与发受链）}]
  \item[\eventanchor{ML-1608-004}{1608（万历三十六年）：矿税、加派、漕运和辽饷的本年奏疏与账目提供财政压力线索，预算、欠额和实解不得混同。（材料核查重点：报称信息的来源、抄传与行政分类）}]
  \item[\eventanchor{ML-1608-005}{1608（万历三十六年）：本年灾害、疫病或地方赈济线索应以受灾地区文书和地方志复核，中央奏报不能覆盖全部社会。（材料核查重点：报称信息的来源、抄传与行政分类）}]
  \item[\eventanchor{ML-1608-006}{1608（万历三十六年）：东林、书院、刻书和宗教活动的本年材料记录精英网络，社团存在不等于一致政治立场。}]
  \item[\eventanchor{ML-1609-001}{1609（万历三十七年）：贾运河竣工}]
  \item[\eventanchor{ML-1609-002}{1609（万历三十七年）：矿税、加派、漕运和辽饷的本年奏疏与账目提供财政压力线索，预算、欠额和实解不得混同。}]
  \item[\eventanchor{ML-1609-003}{1609（万历三十七年）：万历末至天启初的立储、内阁和言路材料标记中枢运作的堵塞与调整，派系归类需回到具体文书。（材料核查重点：诏令或奏议的形成与发受链）}]
  \item[\eventanchor{ML-1609-004}{1609（万历三十七年）：辽东战事、边镇防务和西南兵事的本年军令记录资源调动，战报数字应与朝鲜、满文或地方材料比照。}]
  \item[\eventanchor{ML-1609-005}{1609（万历三十七年）：万历末至天启初的立储、内阁和言路材料标记中枢运作的堵塞与调整，派系归类需回到具体文书。（材料核查重点：报称信息的来源、抄传与行政分类）}]
  \item[\eventanchor{ML-1610-001}{1610（万历三十八年）：李约瑟估计，欧洲文明在天文学和物理学方面大约在这个时候超越了中国}]
  \item[\eventanchor{ML-1610-002}{1610（万历三十八年）：辽东战事、边镇防务和西南兵事的本年军令记录资源调动，战报数字应与朝鲜、满文或地方材料比照。}]
  \item[\eventanchor{ML-1610-003}{1610（万历三十八年）：矿税、加派、漕运和辽饷的本年奏疏与账目提供财政压力线索，预算、欠额和实解不得混同。（材料核查重点：诏令或奏议的形成与发受链）}]
  \item[\eventanchor{ML-1610-004}{1610（万历三十八年）：后金兴起、朝鲜交涉和海上网络的本年多方材料揭示区域关系变化，政权名称与译名须按原文说明。}]
  \item[\eventanchor{ML-1610-005}{1610（万历三十八年）：矿税、加派、漕运和辽饷的本年奏疏与账目提供财政压力线索，预算、欠额和实解不得混同。（材料核查重点：报称信息的来源、抄传与行政分类）}]
  \item[\eventanchor{ML-1610-006}{1610（万历三十八年）：万历末至天启初的立储、内阁和言路材料标记中枢运作的堵塞与调整，派系归类需回到具体文书。}]
  \item[\eventanchor{ML-1611-001}{1611（万历三十九年）：辽东战事、边镇防务和西南兵事的本年军令记录资源调动，战报数字应与朝鲜、满文或地方材料比照。（材料核查重点：诏令或奏议的形成与发受链）}]
  \item[\eventanchor{ML-1611-002}{1611（万历三十九年）：后金兴起、朝鲜交涉和海上网络的本年多方材料揭示区域关系变化，政权名称与译名须按原文说明。}]
  \item[\eventanchor{ML-1611-003}{1611（万历三十九年）：辽东战事、边镇防务和西南兵事的本年军令记录资源调动，战报数字应与朝鲜、满文或地方材料比照。（材料核查重点：报称信息的来源、抄传与行政分类）}]
  \item[\eventanchor{ML-1611-004}{1611（万历三十九年）：东林、书院、刻书和宗教活动的本年材料记录精英网络，社团存在不等于一致政治立场。}]
  \item[\eventanchor{ML-1611-005}{1611（万历三十九年）：辽东战事、边镇防务和西南兵事的本年军令记录资源调动，战报数字应与朝鲜、满文或地方材料比照。（材料核查重点：同年地方材料所见的执行范围）}]
  \item[\eventanchor{ML-1611-006}{1611（万历三十九年）：矿税、加派、漕运和辽饷的本年奏疏与账目提供财政压力线索，预算、欠额和实解不得混同。}]
  \item[\eventanchor{ML-1612-001}{1612（万历四十年）：后金兴起、朝鲜交涉和海上网络的本年多方材料揭示区域关系变化，政权名称与译名须按原文说明。（材料核查重点：诏令或奏议的形成与发受链）}]
  \item[\eventanchor{ML-1612-002}{1612（万历四十年）：东林、书院、刻书和宗教活动的本年材料记录精英网络，社团存在不等于一致政治立场。}]
  \item[\eventanchor{ML-1612-003}{1612（万历四十年）：后金兴起、朝鲜交涉和海上网络的本年多方材料揭示区域关系变化，政权名称与译名须按原文说明。（材料核查重点：报称信息的来源、抄传与行政分类）}]
  \item[\eventanchor{ML-1612-004}{1612（万历四十年）：本年灾害、疫病或地方赈济线索应以受灾地区文书和地方志复核，中央奏报不能覆盖全部社会。}]
  \item[\eventanchor{ML-1612-005}{1612（万历四十年）：后金兴起、朝鲜交涉和海上网络的本年多方材料揭示区域关系变化，政权名称与译名须按原文说明。（材料核查重点：同年地方材料所见的执行范围）}]
  \item[\eventanchor{ML-1613-001}{1613（万历四十一年）：东林、书院、刻书和宗教活动的本年材料记录精英网络，社团存在不等于一致政治立场。（材料核查重点：诏令或奏议的形成与发受链）}]
  \item[\eventanchor{ML-1613-002}{1613（万历四十一年）：本年灾害、疫病或地方赈济线索应以受灾地区文书和地方志复核，中央奏报不能覆盖全部社会。}]
  \item[\eventanchor{ML-1613-003}{1613（万历四十一年）：东林、书院、刻书和宗教活动的本年材料记录精英网络，社团存在不等于一致政治立场。（材料核查重点：报称信息的来源、抄传与行政分类）}]
  \item[\eventanchor{ML-1613-004}{1613（万历四十一年）：万历末至天启初的立储、内阁和言路材料标记中枢运作的堵塞与调整，派系归类需回到具体文书。}]
  \item[\eventanchor{ML-1613-005}{1613（万历四十一年）：东林、书院、刻书和宗教活动的本年材料记录精英网络，社团存在不等于一致政治立场。（材料核查重点：同年地方材料所见的执行范围）}]
  \item[\eventanchor{ML-1613-006}{1613（万历四十一年）：后金兴起、朝鲜交涉和海上网络的本年多方材料揭示区域关系变化，政权名称与译名须按原文说明。}]
  \item[\eventanchor{ML-1614-001}{1614（万历四十二年）：本年灾害、疫病或地方赈济线索应以受灾地区文书和地方志复核，中央奏报不能覆盖全部社会。（材料核查重点：诏令或奏议的形成与发受链）}]
  \item[\eventanchor{ML-1614-002}{1614（万历四十二年）：万历末至天启初的立储、内阁和言路材料标记中枢运作的堵塞与调整，派系归类需回到具体文书。}]
  \item[\eventanchor{ML-1614-003}{1614（万历四十二年）：本年灾害、疫病或地方赈济线索应以受灾地区文书和地方志复核，中央奏报不能覆盖全部社会。（材料核查重点：报称信息的来源、抄传与行政分类）}]
  \item[\eventanchor{ML-1614-004}{1614（万历四十二年）：矿税、加派、漕运和辽饷的本年奏疏与账目提供财政压力线索，预算、欠额和实解不得混同。}]
  \item[\eventanchor{ML-1614-005}{1614（万历四十二年）：本年灾害、疫病或地方赈济线索应以受灾地区文书和地方志复核，中央奏报不能覆盖全部社会。（材料核查重点：同年地方材料所见的执行范围）}]
  \item[\eventanchor{ML-1614-006}{1614（万历四十二年）：东林、书院、刻书和宗教活动的本年材料记录精英网络，社团存在不等于一致政治立场。}]
  \item[\eventanchor{ML-1615-001}{1615（万历四十三年）：努尔哈赤派遣最后一位朝贡使者前往北京}]
  \item[\eventanchor{ML-1615-002}{1615（万历四十三年）：矿税、加派、漕运和辽饷的本年奏疏与账目提供财政压力线索，预算、欠额和实解不得混同。}]
  \item[\eventanchor{ML-1615-003}{1615（万历四十三年）：万历末至天启初的立储、内阁和言路材料标记中枢运作的堵塞与调整，派系归类需回到具体文书。（材料核查重点：诏令或奏议的形成与发受链）}]
  \item[\eventanchor{ML-1615-004}{1615（万历四十三年）：辽东战事、边镇防务和西南兵事的本年军令记录资源调动，战报数字应与朝鲜、满文或地方材料比照。}]
  \item[\eventanchor{ML-1615-005}{1615（万历四十三年）：万历末至天启初的立储、内阁和言路材料标记中枢运作的堵塞与调整，派系归类需回到具体文书。（材料核查重点：报称信息的来源、抄传与行政分类）}]
  \item[\eventanchor{ML-1615-006}{1615（万历四十三年）：本年灾害、疫病或地方赈济线索应以受灾地区文书和地方志复核，中央奏报不能覆盖全部社会。}]
  \item[\eventanchor{ML-1616-002}{1616（万历四十四年）：辽东战事、边镇防务和西南兵事的本年军令记录资源调动，战报数字应与朝鲜、满文或地方材料比照。}]
  \item[\eventanchor{ML-1616-003}{1616（万历四十四年）：矿税、加派、漕运和辽饷的本年奏疏与账目提供财政压力线索，预算、欠额和实解不得混同。（材料核查重点：诏令或奏议的形成与发受链）}]
  \item[\eventanchor{ML-1616-004}{1616（万历四十四年）：后金兴起、朝鲜交涉和海上网络的本年多方材料揭示区域关系变化，政权名称与译名须按原文说明。}]
  \item[\eventanchor{ML-1616-005}{1616（万历四十四年）：矿税、加派、漕运和辽饷的本年奏疏与账目提供财政压力线索，预算、欠额和实解不得混同。（材料核查重点：报称信息的来源、抄传与行政分类）}]
  \item[\eventanchor{ML-1616-006}{1616（万历四十四年）：万历末至天启初的立储、内阁和言路材料标记中枢运作的堵塞与调整，派系归类需回到具体文书。}]
  \item[\eventanchor{ML-1617-001}{1617（万历四十五年）：万历末至天启初的立储、内阁和言路材料标记中枢运作的堵塞与调整，派系归类需回到具体文书。}]
  \item[\eventanchor{ML-1617-002}{1617（万历四十五年）：后金兴起、朝鲜交涉和海上网络的本年多方材料揭示区域关系变化，政权名称与译名须按原文说明。}]
  \item[\eventanchor{ML-1617-003}{1617（万历四十五年）：辽东战事、边镇防务和西南兵事的本年军令记录资源调动，战报数字应与朝鲜、满文或地方材料比照。（材料核查重点：诏令或奏议的形成与发受链）}]
  \item[\eventanchor{ML-1617-004}{1617（万历四十五年）：东林、书院、刻书和宗教活动的本年材料记录精英网络，社团存在不等于一致政治立场。}]
  \item[\eventanchor{ML-1617-005}{1617（万历四十五年）：辽东战事、边镇防务和西南兵事的本年军令记录资源调动，战报数字应与朝鲜、满文或地方材料比照。（材料核查重点：报称信息的来源、抄传与行政分类）}]
  \item[\eventanchor{ML-1617-006}{1617（万历四十五年）：矿税、加派、漕运和辽饷的本年奏疏与账目提供财政压力线索，预算、欠额和实解不得混同。}]
  \item[\eventanchor{ML-1618-001}{1618（万历四十六年）五月9日：抚顺之战：后金占领抚顺}]
  \item[\eventanchor{ML-1618-002}{1618（万历四十六年）夏：清河之战：后金攻克清河}]
  \item[\eventanchor{ML-1618-003}{1618（万历四十六年）：后金兴起、朝鲜交涉和海上网络的本年多方材料揭示区域关系变化，政权名称与译名须按原文说明。（材料核查重点：诏令或奏议的形成与发受链）}]
  \item[\eventanchor{ML-1618-004}{1618（万历四十六年）：本年灾害、疫病或地方赈济线索应以受灾地区文书和地方志复核，中央奏报不能覆盖全部社会。}]
  \item[\eventanchor{ML-1618-005}{1618（万历四十六年）：后金兴起、朝鲜交涉和海上网络的本年多方材料揭示区域关系变化，政权名称与译名须按原文说明。（材料核查重点：报称信息的来源、抄传与行政分类）}]
  \item[\eventanchor{ML-1619-001}{1619（万历四十七年）四月18日：萨尔浒之战：明军被后金歼灭}]
  \item[\eventanchor{ML-1619-003}{1619（万历四十七年）九月3日：铁岭之战：后金占领铁岭}]
  \item[\eventanchor{ML-1619-004}{1619（万历四十七年）九月：首位俄罗斯特使伊万·佩特林抵达北京}]
  \item[\eventanchor{ML-1619-005}{1619（万历四十七年）：东林、书院、刻书和宗教活动的本年材料记录精英网络，社团存在不等于一致政治立场。}]
  \item[\eventanchor{ML-1620-001}{1620（万历四十八年）八月18日：万历皇帝驾崩}]
  \item[\eventanchor{ML-1620-002}{1620（万历四十八年）八月28日：朱常洛即位太常皇帝}]
  \item[\eventanchor{ML-1620-003}{1620（万历四十八年）九月6日：太常皇帝病了}]
  \item[\eventanchor{ML-1620-004}{1620（万历四十八年）九月26日：太常帝驾崩}]
  \item[\eventanchor{ML-1620-005}{1620（万历四十八年）十月1日：朱由娇成为天启皇帝}]
  \item[\eventanchor{ML-1620-006}{1620（万历四十八年）：明代铸造厂开始生产红衣袍。}]
\end{description}
